\chapter{Derivation of the limiting kernel on \(S^1\)}
\label{app:continuum-kernel}

This appendix derives the explicit form of the limiting neural tangent kernel on the circle
for the model introduced in Chapter~\ref{cha:methodology}, and computes the corresponding
Fourier eigenvalues of the integral operator \(T\). These derivations justify
Proposition~\ref{prop:results-continuum-ntk} and
Theorem~\ref{thm:results-continuum-spectrum}.

\section{Decomposition of the empirical neural tangent kernel}
\label{sec:app-ntk-decomposition}

Recall the network
\[
	f(x;\theta)=\frac{1}{\sqrt m}\sum_{i=1}^m a_i \sigma(w_i^\top x+b_i),
\]
where \(x\in S^1\subset\mathbb R^2\), \(\sigma(t)=t_+\), and
\(\theta=\{(a_i,w_i,b_i)\}_{i=1}^m\). For \(x,y\in S^1\), the empirical neural tangent kernel is
\begin{equation}
	\Theta_m(x,y)
	=
	\bigl\langle \nabla_\theta f(x;\theta),\nabla_\theta f(y;\theta)\bigr\rangle.
	\label{eq:app-empirical-ntk}
\end{equation}

For each hidden unit \(i\), define
\[
	u_i:=w_i^\top x+b_i,
	\qquad
	v_i:=w_i^\top y+b_i.
\]
Then
\[
	\frac{\partial f(x;\theta)}{\partial a_i}
	=
	\frac{1}{\sqrt m}\sigma(u_i),
	\qquad
	\frac{\partial f(x;\theta)}{\partial w_i}
	=
	\frac{1}{\sqrt m}a_i\sigma'(u_i)\,x,
	\qquad
	\frac{\partial f(x;\theta)}{\partial b_i}
	=
	\frac{1}{\sqrt m}a_i\sigma'(u_i).
\]
Substituting into \eqref{eq:app-empirical-ntk} yields
\begin{equation}
	\Theta_m(x,y)
	=
	\Theta_m^{(a)}(x,y)+\Theta_m^{(w)}(x,y)+\Theta_m^{(b)}(x,y),
	\label{eq:app-ntk-decomposition}
\end{equation}
where
\begin{align*}
	\Theta_m^{(a)}(x,y)
	 & =
	\frac1m\sum_{i=1}^m \sigma(u_i)\sigma(v_i),                   \\
	\Theta_m^{(w)}(x,y)
	 & =
	\frac1m\sum_{i=1}^m a_i^2 \sigma'(u_i)\sigma'(v_i)\,x^\top y, \\
	\Theta_m^{(b)}(x,y)
	 & =
	\frac1m\sum_{i=1}^m a_i^2 \sigma'(u_i)\sigma'(v_i).
\end{align*}

\section{Infinite-width limit}
\label{sec:app-ntk-limit}

At initialization, the parameters satisfy
\[
	a_i\sim \mathcal N(0,1),
	\qquad
	w_i\sim \mathcal N(0,I_2),
	\qquad
	b_i(0)=0,
\]
independently across \(i\). Since the hidden units are independent and identically distributed,
the law of large numbers implies that \(\Theta_m(x,y)\) converges almost surely, as
\(m\to\infty\), to a deterministic kernel \(\Theta(x,y)\).

Because \(b_i(0)=0\), the limiting kernel is
\begin{equation}
	\Theta(x,y)
	=
	\mathbb E\bigl[\sigma(u)\sigma(v)\bigr]
	+
	\bigl(x^\top y+1\bigr)\mathbb E\bigl[\sigma'(u)\sigma'(v)\bigr],
	\label{eq:app-limit-kernel-general}
\end{equation}
where
\[
	u=w^\top x,
	\qquad
	v=w^\top y,
	\qquad
	w\sim \mathcal N(0,I_2).
\]
We write
\[
	K_0(x,y):=\mathbb E\bigl[\sigma(u)\sigma(v)\bigr],
	\qquad
	K_1(x,y):=\mathbb E\bigl[\sigma'(u)\sigma'(v)\bigr],
\]
so that
\[
	\Theta(x,y)=K_0(x,y)+\bigl(x^\top y+1\bigr)K_1(x,y).
\]

\section{Restriction to the circle}
\label{sec:app-kernel-on-circle}

Write
\[
	x(\varphi)=(\cos\varphi,\sin\varphi),
	\qquad
	x(\psi)=(\cos\psi,\sin\psi),
\]
and let \(\delta\in[0,\pi]\) denote the angle between \(x(\varphi)\) and \(x(\psi)\), so that
\[
	x(\varphi)^\top x(\psi)=\cos\delta.
\]
By rotational invariance of the Gaussian measure, the kernel depends only on \(\delta\). We
therefore write \(\Theta(\delta)\), \(K_0(\delta)\), and \(K_1(\delta)\).

To compute these quantities, we rotate coordinates so that
\[
	x=(1,0),
	\qquad
	y=(\cos\delta,\sin\delta).
\]
If \(w=(Z_1,Z_2)\) with \(Z_1,Z_2\overset{\mathrm{iid}}{\sim}\mathcal N(0,1)\), then
\[
	u=Z_1,
	\qquad
	v=Z_1\cos\delta+Z_2\sin\delta.
\]

\begin{lemma}
	\label{lem:app-K1}
	For \(\delta\in[0,\pi]\),
	\begin{equation}
		K_1(\delta)=\mathbb P(u>0,v>0)=\frac{\pi-\delta}{2\pi}.
		\label{eq:app-K1-explicit}
	\end{equation}
\end{lemma}

\begin{proof}
	Since \((u,v)\) is a centered bivariate Gaussian with correlation \(\cos\delta\), the
	probability that both coordinates are positive is the standard Gaussian quadrant probability,
	which equals \((\pi-\delta)/(2\pi)\).
\end{proof}

\begin{lemma}
	\label{lem:app-K0}
	For \(\delta\in[0,\pi]\),
	\begin{equation}
		K_0(\delta)=\mathbb E[\sigma(u)\sigma(v)]
		=
		\frac{1}{2\pi}\Bigl(\sin\delta+(\pi-\delta)\cos\delta\Bigr).
		\label{eq:app-K0-explicit}
	\end{equation}
\end{lemma}

\begin{proof}
	This is the standard arc-cosine covariance formula for ReLU features in two dimensions.
\end{proof}

\begin{proof}[Proof of Proposition~\ref{prop:results-continuum-ntk}]
	Using \eqref{eq:app-limit-kernel-general}, \eqref{eq:app-K1-explicit}, and
	\eqref{eq:app-K0-explicit}, we obtain
	\begin{align}
		\Theta(\delta)
		 & =
		\frac{1}{2\pi}\Bigl(\sin\delta+(\pi-\delta)\cos\delta\Bigr)
		+
		\frac{1+\cos\delta}{2\pi}(\pi-\delta) \notag \\
		 & =
		\frac{1}{2\pi}
		\Bigl(
		\sin\delta+2(\pi-\delta)\cos\delta+(\pi-\delta)
		\Bigr).
		\label{eq:app-Theta-explicit}
	\end{align}
	This is exactly \eqref{eq:results-continuum-ntk}. The convolution form of the operator
	follows from the fact that \(\Theta\) depends only on the angular difference.
\end{proof}

\section{Fourier coefficients of the limiting kernel}
\label{sec:app-fourier-coefficients}

Because the operator \(T\) is a convolution operator, it is diagonalized by the real Fourier
basis. If \(\lambda_k\) denotes the eigenvalue associated with frequency \(k\), then
\begin{equation}
	\lambda_k
	=
	\frac{1}{\pi}\int_0^\pi \Theta(\delta)\cos(k\delta)\,d\delta,
	\qquad
	k\ge 0.
	\label{eq:app-lambda-k}
\end{equation}
Substituting \eqref{eq:app-Theta-explicit} gives
\[
	\lambda_k
	=
	\frac{1}{2\pi^2}
	\int_0^\pi
	\Bigl(
	\sin\delta+2(\pi-\delta)\cos\delta+(\pi-\delta)
	\Bigr)\cos(k\delta)\,d\delta.
\]

We first record a basic integration lemma.

\begin{lemma}
	\label{lem:app-Im}
	For every integer \(m\ge 1\),
	\begin{equation}
		I_m:=\int_0^\pi (\pi-\delta)\cos(m\delta)\,d\delta
		=
		\frac{1-(-1)^m}{m^2}.
		\label{eq:app-Im}
	\end{equation}
\end{lemma}

\begin{proof}
	Integrating by parts gives
	\[
		I_m
		=
		\left[(\pi-\delta)\frac{\sin(m\delta)}{m}\right]_0^\pi
		+
		\frac1m\int_0^\pi \sin(m\delta)\,d\delta.
	\]
	The boundary term vanishes, and
	\[
		\frac1m\int_0^\pi \sin(m\delta)\,d\delta
		=
		\frac1m\left[-\frac{\cos(m\delta)}{m}\right]_0^\pi
		=
		\frac{1-(-1)^m}{m^2}.
	\]
\end{proof}

\begin{proof}[Proof of Theorem~\ref{thm:results-continuum-spectrum}]
	For \(k\ge 2\), write
	\[
		\lambda_k=\frac{1}{2\pi^2}(A_k+B_k+C_k),
	\]
	where
	\begin{align*}
		A_k & :=\int_0^\pi \sin\delta\,\cos(k\delta)\,d\delta,            \\
		B_k & :=\int_0^\pi 2(\pi-\delta)\cos\delta\cos(k\delta)\,d\delta, \\
		C_k & :=\int_0^\pi (\pi-\delta)\cos(k\delta)\,d\delta.
	\end{align*}

	Using
	\[
		\sin\delta\cos(k\delta)
		=
		\frac12\bigl(\sin((k+1)\delta)-\sin((k-1)\delta)\bigr),
	\]
	we obtain
	\[
		A_k
		=
		\frac12\left(
		\frac{1-(-1)^{k+1}}{k+1}
		-
		\frac{1-(-1)^{k-1}}{k-1}
		\right).
	\]
	Using
	\[
		2\cos\delta\cos(k\delta)=\cos((k+1)\delta)+\cos((k-1)\delta),
	\]
	together with Lemma~\ref{lem:app-Im}, we obtain
	\[
		B_k
		=
		\frac{1-(-1)^{k+1}}{(k+1)^2}
		+
		\frac{1-(-1)^{k-1}}{(k-1)^2},
		\qquad
		C_k
		=
		\frac{1-(-1)^k}{k^2}.
	\]

	If \(k\ge 3\) is odd, then \(k\pm1\) are even, so \(A_k=B_k=0\), while
	\(C_k=2/k^2\). Hence
	\begin{equation}
		\lambda_k=\frac{1}{\pi^2 k^2},
		\qquad
		k\ge 3 \text{ odd}.
		\label{eq:app-lambda-odd}
	\end{equation}

	If \(k\ge 2\) is even, then \(k\pm1\) are odd, so \(C_k=0\), and
	\[
		A_k
		=
		-\frac{2}{k^2-1},
		\qquad
		B_k
		=
		\frac{4(k^2+1)}{(k^2-1)^2}.
	\]
	Therefore
	\[
		A_k+B_k
		=
		\frac{2(k^2+3)}{(k^2-1)^2},
	\]
	and so
	\begin{equation}
		\lambda_k
		=
		\frac{k^2+3}{\pi^2 (k^2-1)^2},
		\qquad
		k\ge 2 \text{ even}.
		\label{eq:app-lambda-even}
	\end{equation}

	It remains to compute \(k=0\) and \(k=1\) directly. For \(k=0\),
	\[
		\lambda_0
		=
		\frac{1}{2\pi^2}
		\left(
		\int_0^\pi \sin\delta\,d\delta
		+
		2\int_0^\pi (\pi-\delta)\cos\delta\,d\delta
		+
		\int_0^\pi (\pi-\delta)\,d\delta
		\right)
		=
		\frac14+\frac{3}{\pi^2}.
	\]
	For \(k=1\),
	\[
		\lambda_1
		=
		\frac{1}{2\pi^2}
		\left(
		\int_0^\pi \sin\delta\cos\delta\,d\delta
		+
		2\int_0^\pi (\pi-\delta)\cos^2\delta\,d\delta
		+
		\int_0^\pi (\pi-\delta)\cos\delta\,d\delta
		\right)
		=
		\frac14+\frac{1}{\pi^2}.
	\]
	This proves \eqref{eq:results-lambda-0-1} and \eqref{eq:results-lambda-k}.
\end{proof}

\section{Effect of the bias contribution}
\label{sec:app-bias-contribution}

We now isolate the effect of the trainable hidden bias. From the NTK decomposition in
\eqref{eq:app-ntk-decomposition}, the limiting kernel can be written as
\[
	\Theta(x,y)
	=
	K_0(x,y)
	+
	x^\top y\,K_1(x,y)
	+
	K_1(x,y).
\]
The first two terms come from differentiating with respect to the output weights and input weights.
The last term comes from differentiating with respect to the hidden biases. Thus, if the trainable
hidden-bias contribution is removed, the limiting kernel becomes
\[
	\Theta^{(\mathrm{nb})}(x,y)
	=
	K_0(x,y)+x^\top y\,K_1(x,y).
\]
On the circle, \(x^\top y=\cos\delta\), so by \eqref{eq:app-K0-explicit} and
\eqref{eq:app-K1-explicit},
\begin{align}
	\Theta^{(\mathrm{nb})}(\delta)
	 & =
	\frac{1}{2\pi}
	\Bigl(
	\sin\delta+(\pi-\delta)\cos\delta
	\Bigr)
	+
	\cos\delta\,\frac{\pi-\delta}{2\pi} \notag \\
	 & =
	\frac{1}{2\pi}
	\Bigl(
	\sin\delta+2(\pi-\delta)\cos\delta
	\Bigr).
	\label{eq:app-no-bias-kernel}
\end{align}

The Fourier eigenvalues of this no-bias kernel are
\[
	\lambda_k^{(\mathrm{nb})}
	=
	\frac{1}{\pi}
	\int_0^\pi
	\Theta^{(\mathrm{nb})}(\delta)\cos(k\delta)\,d\delta.
\]
Using \eqref{eq:app-no-bias-kernel}, this becomes
\[
	\lambda_k^{(\mathrm{nb})}
	=
	\frac{1}{2\pi^2}
	\int_0^\pi
	\Bigl(
	\sin\delta+2(\pi-\delta)\cos\delta
	\Bigr)
	\cos(k\delta)\,d\delta.
\]
In the notation of the proof of Theorem~\ref{thm:results-continuum-spectrum},
\[
	\lambda_k^{(\mathrm{nb})}
	=
	\frac{1}{2\pi^2}(A_k+B_k),
\]
whereas the full eigenvalue is
\[
	\lambda_k
	=
	\frac{1}{2\pi^2}(A_k+B_k+C_k).
\]
Thus the term \(C_k\), which comes from
\[
	\frac{1}{2\pi}(\pi-\delta),
\]
is precisely the trainable-bias contribution.

For \(k=0\), direct integration gives
\[
	\lambda_0^{(\mathrm{nb})}
	=
	\frac{1}{2\pi^2}
	\left(
	\int_0^\pi \sin\delta\,d\delta
	+
	2\int_0^\pi(\pi-\delta)\cos\delta\,d\delta
	\right).
\]
Since
\[
	\int_0^\pi \sin\delta\,d\delta=2,
	\qquad
	\int_0^\pi(\pi-\delta)\cos\delta\,d\delta=2,
\]
we obtain
\[
	\lambda_0^{(\mathrm{nb})}
	=
	\frac{1}{2\pi^2}(2+4)
	=
	\frac{3}{\pi^2}.
\]

For \(k=1\),
\[
	\lambda_1^{(\mathrm{nb})}
	=
	\frac{1}{2\pi^2}
	\left(
	\int_0^\pi \sin\delta\cos\delta\,d\delta
	+
	2\int_0^\pi(\pi-\delta)\cos^2\delta\,d\delta
	\right).
\]
The first integral is zero. For the second, using
\[
	\cos^2\delta=\frac{1+\cos(2\delta)}{2},
\]
we get
\[
	2\int_0^\pi(\pi-\delta)\cos^2\delta\,d\delta
	=
	\int_0^\pi(\pi-\delta)\,d\delta
	+
	\int_0^\pi(\pi-\delta)\cos(2\delta)\,d\delta.
\]
The second term is zero by Lemma~\ref{lem:app-Im}, since \(2\) is even. Hence
\[
	2\int_0^\pi(\pi-\delta)\cos^2\delta\,d\delta
	=
	\frac{\pi^2}{2},
\]
and therefore
\[
	\lambda_1^{(\mathrm{nb})}
	=
	\frac14.
\]

For \(k\ge 2\), we use the values of \(A_k\) and \(B_k\) computed in the proof of
Theorem~\ref{thm:results-continuum-spectrum}. If \(k\ge 3\) is odd, then \(k-1\) and \(k+1\)
are even, so
\[
	A_k=0,
	\qquad
	B_k=0.
\]
Therefore
\begin{equation}
	\lambda_k^{(\mathrm{nb})}=0,
	\qquad
	k\ge 3 \text{ odd}.
	\label{eq:app-no-bias-odd}
\end{equation}
If \(k\ge 2\) is even, then
\[
	A_k=-\frac{2}{k^2-1},
	\qquad
	B_k=\frac{4(k^2+1)}{(k^2-1)^2}.
\]
Thus
\[
	A_k+B_k
	=
	\frac{2(k^2+3)}{(k^2-1)^2},
\]
and hence
\begin{equation}
	\lambda_k^{(\mathrm{nb})}
	=
	\frac{k^2+3}{\pi^2(k^2-1)^2},
	\qquad
	k\ge 2 \text{ even}.
	\label{eq:app-no-bias-even}
\end{equation}

This proves Corollary~\ref{cor:results-bias-spectrum}. Comparing the no-bias eigenvalues with
the full eigenvalues in Theorem~\ref{thm:results-continuum-spectrum}, the trainable-bias
contribution is
\[
	\lambda_0-\lambda_0^{(\mathrm{nb})}
	=
	\frac14,
	\qquad
	\lambda_1-\lambda_1^{(\mathrm{nb})}
	=
	\frac{1}{\pi^2},
\]
while for \(k\ge 2\),
\[
	\lambda_k-\lambda_k^{(\mathrm{nb})}
	=
	\begin{cases}
		0,                   & k\ge 2 \text{ even}, \\[0.8em]
		\dfrac{1}{\pi^2k^2}, & k\ge 3 \text{ odd}.
	\end{cases}
\]
Thus the bias term does not change the even frequencies \(k\ge 2\), but it is exactly what makes
the odd frequencies \(k\ge 3\) visible in the fixed-kernel dynamics. The frequency \(k=1\) is
exceptional: it already has eigenvalue \(1/4\) without the bias, and the bias adds the additional
\(1/\pi^2\) contribution.
